\chapter{KaiMing Gate：TMM 风险、验收与论文写作前置条件}
\label{chap:risks}

本文的主要结论仍然保守，但主证据已比上一版更强。P0 E2 表明 TRR 在当前 guitar-effects、204-query、retrieval-only benchmark 上，是 Wav2Vec、FeatureNN、CLAP、PaSST 和 PANNs 中参数对齐最强的 retrieval prior；P0 E3 表明真实音频退化下，adaptive fusion 明显优于 audio-heavy fixed fusion，并接近中高文本权重的稳定工作区间。该结论足以构成一篇围绕 retrieval-grounded editable audio effect control 与 real-degradation fusion robustness 的 TMM 方法与实验证据报告，但不足以支撑通用音频表示、端到端主观质量或完整自动化音频系统的强 claim。



KaiMing gate 用于防止 TMM 论文在证据尚未闭合时提前固化。当前阶段允许修订 abstract、method skeleton、experiment checklist、supplement 和 response letter；不允许锁定强 novelty claim、SOTA-level conclusion 或未绑定 artifact 的 reviewer response。主论文启动条件应由表~\ref{tab:kaiming-gates} 控制。

\begin{table}[H]
	\centering
	\caption{KaiMing gate 与 TMM 论文写作前置条件。}
	\label{tab:kaiming-gates}
	\scriptsize
	\resizebox{\textwidth}{!}{%
		\begin{tabular}{>{\raggedright\arraybackslash}p{0.17\textwidth}>{\raggedright\arraybackslash}p{0.31\textwidth}>{\raggedright\arraybackslash}p{0.26\textwidth}>{\raggedright\arraybackslash}p{0.18\textwidth}}
		\toprule
		Gate & 必须通过的证据 & 失败时的论文处理 & 状态 \\
		\midrule
			范围门 & TMM scope、主问题、输出对象和贡献边界已冻结 & 回到引言与相关工作，不写泛化 agent claim & 正在迁移。 \\
				基线门 & CLAP、PaSST、PANNs 已纳入；直接回归和 Text2FX 式基线状态被标注 & 主表不能声称通用 SOTA 级优势 & 部分完成。 \\
			消融门 & 二阶纹理、层选择、归一化和 top-k 有可解释贡献 & 方法贡献降级为经验检索特征 & 待实验。 \\
			泄漏门 & exact duplicate、near duplicate 和 family split 风险已审计 & 降低 claim 或更换主协议 & 部分完成。 \\
				鲁棒门 & P0 E3 已覆盖 AWGN、MP3、reverb、truncate；conflict 场景仍属旧诊断 & 鲁棒性写为当前退化集合下的 fusion 机制证据 & 部分完成。 \\
			感知门 & 参数指标与听感评分关系已估计，或明确相关性不足 & 参数指标限定为 controllability proxy & 待补强。 \\
			复现门 & logger、checksum、run scripts、匿名包和 checklist 对齐 & 只能作为内部报告或草稿 & 待实现。 \\
			提交门 & 主文、附录、答复、checklist 和 LaTeX warning 通过 TMM/SPS 检查 & 不进入正式提交 & 待全量检查。 \\
		\bottomrule
	\end{tabular}}
\end{table}

第一个局限是数据域。当前 benchmark 主要围绕吉他效果预设展开，参数结构也绑定于特定效果链。TRR 对该域有效，并不自动意味着它对鼓组、人声、混音母带或开放世界音频编辑同样有效。第二个局限是指标。Param. Dist.、Acc@0.1、Recall、Cosine 和 Module 都是参数空间指标，不能完全替代听感距离。第三个局限是协议。严格 split 已消除 exact resolved-audio reuse，但 near-duplicate 风险仍需更强审计。第四个局限是工程证据。产品链路存在，但主实验没有直接评价完整插件输出音频。

第二个需要明确的风险是融合边界。P0 E3 说明 adaptive fusion 在 AWGN、MP3、reverb 和 truncate 退化下接近稳定工作区间，但它不是每个 condition / 每个 metric 的压倒性赢家。旧 Protocol-C conflict 场景中的 Fusion L2=4.5259 远高于 TRR-only 的 0.2970，仍说明当前融合规则无法可靠处理相互矛盾的文本与音频输入。若后续要扩展融合主张，必须新增独立协议，定义冲突输入、真实退化、失败类型和统计检验。

\begin{table}[H]
	\centering
	\caption{主张、证据与允许措辞矩阵。}
	\label{tab:claim-matrix}
	\small
	\resizebox{\textwidth}{!}{%
	\begin{tabular}{lll}
		\toprule
			主张 & 证据状态 & 允许措辞 \\
		\midrule
			TRR 优于当前已评估检索基线 & P0 E2 支撑 & 可以作为主结论。 \\
		TRR 是通用最优音频表示 & 无跨域证据 & 不应声称。 \\
			Adaptive fusion 具备真实退化稳健性 & P0 E3 支撑 & 可作为主结论的第二部分。 \\
			Adaptive fusion 全面优于所有固定策略 & per-condition 表不支持 & 不应声称。 \\
		TRR-based system 有主观感知价值 & 听测补充支持 & 可保守表述为主观背景。 \\
		系统严格优于人工调参或 MusicGen & 元数据与统计边界不足 & 不应声称。 \\
			P0 PaSST/PANNs 已正式验证 & 用户已 double check 真实完整正确 & 可以进入论文主结果。 \\
		\bottomrule
	\end{tabular}}
\end{table}

表~\ref{tab:claim-matrix} 是本文最终结论的约束条件。后续交付应按三个优先级推进。首先，冻结主证据：将 P0 E2、P0 E3 和 listening test 作为 canonical evidence，并明确各自的统计口径、数据集快照和局限边界。其次，收束 TMM 论文包：主文、supplement、response letter 和 revision checklist 必须引用同一组证据，并清理页数、warning、citation 与匿名包问题。最后，若要主张跨域鲁棒性或端到端听感优势，还需要新增独立协议。

在当前证据条件下，最合适的论文式结论是：TRR 将深层音频特征的二阶纹理统计用于可执行音频效果参数检索，在当前 guitar-effects 检索基准中，相比 Wav2Vec、FeatureNN、CLAP、PaSST 和 PANNs 等已评估基线表现出更好的参数对齐；同时，真实退化下的 adaptive fusion 比 audio-heavy fixed fusion 更稳健，并接近中高文本权重的稳定工作区间。该结论比上一版更强，但仍窄而可审计。



因此，本报告对下一阶段的监督结论是：先完成 canonical evidence freeze 和 submission hygiene，再决定是否投入新的 direct regression、reranker 或跨域消融实验。当表~\ref{tab:kaiming-gates} 至少通过 scope、baseline honesty、leakage、reproducibility 和 submission hygiene 五项后，才应进入 TMM 主文定稿。若短期内无法通过 ablation 或跨域 gate，项目仍可以形成严谨的 TMM revision package，但主张必须保持为“当前已评估检索基线下的 texture-aware retrieval prior + 当前真实退化集合下的 adaptive fusion robustness”，不得升级为通用 SOTA 或完整系统胜负结论。
